% Restyle of fig-sdt to the tikz-diagram-craft v2 standard (v1 names and raw
% colours migrated to roles; the two densities and pgfplots axis are
% unchanged, ch4.md 4.7). Fixes label crowding: the curve name tags move
% above the peaks with clearance instead of sitting beside them.
%
% Reader question: how does the forced-zoom threshold trade misses against
% false alarms?
% Claim: moving a threshold on fixed evidence distributions changes misses
% and false alarms, not sensitivity. The drawn threshold is illustrative.
% Evidence: Def.~sdt; costs and prevalence determine the preferred threshold.
% Vertical scale is arbitrary; common sigma=sqrt(.7), mean separation=3.
% Must distinguish: the miss region (dangerous mass left of beta) from the
% false-alarm region (safe mass right of beta); which curve is dashed.
% Grammar chosen: two-distribution decision plot (pgfplots) -- axes, named
% densities, a criterion line, and the two error regions in the data field.
% Grammar rejected: text over the curves (the v1 draft) crowds the peaks;
% callouts inside the peaks would sit on the densest ink in the figure.
% Five-second test: which side of beta is the costly error?
\begin{figure}[H]
\centering
\pdfigurehue{pdteal}%
\begin{tikzpicture}[pd figure,
  declare function={
    safe(\x)   = 5*exp(-((\x-3)^2)/1.4);
    danger(\x) = 5*exp(-((\x-6)^2)/1.4);
  }]
\begin{axis}[
  width=7.6cm,height=5.6cm,
  axis lines=left,
  axis line style={pd rule},
  xmin=0,xmax=10,ymin=0,ymax=8.6,
  xtick={0,2,4,6,8,10},ytick={0},
  tick label style={font=\pdfiglabel,text=pdinkmuted},
  xlabel={evidence of danger the operator reads},
  ylabel={relative density},
  xlabel style={align=center,text width=6.8cm},
  every axis label/.style={font=\pdfiglabel,text=pdinkmuted},
  clip=false,
  axis on top,
  ]
  % miss area: dangerous mass left of beta
  \addplot[pd breach fill,domain=0:5.0,samples=80,smooth] {danger(x)} \closedcycle;
  % false-alarm area: safe mass right of beta
  \addplot[pd neutral fill,domain=5.0:10,samples=80,smooth] {safe(x)} \closedcycle;

  \addplot[pd focus rule,line width=1.1pt,domain=0:10,samples=120,smooth] {safe(x)};
  \addplot[pd rule,dashed,line width=1.0pt,domain=0:10,samples=120,smooth] {danger(x)};

  % error labels, low, inside their own tinted region, clear of the axis
  \node[pd label,anchor=south,align=center] at (axis cs:3.6,0.55) {miss};
  \node[pd focus label,anchor=south,align=center] at (axis cs:7.3,0.55) {false alarm};

  % curve name tags, well above each peak, clear of the density's own ink
  \node[pd focus label,anchor=south,align=center,text width=2.2cm] at (axis cs:3,5.85) {safe actions};
  \node[pd label,anchor=south west,text width=2.2cm] at (axis cs:5.55,5.85) {dangerous};

  % criterion beta: full-height rule with its own label in the low valley
  \draw[pd rule,line width=1.0pt] (axis cs:5.0,0) -- (axis cs:5.0,6.7);
  \node[pd label,anchor=north west,text width=1.7cm] at (axis cs:7.5,4.6) {criterion $\beta=5$};

  % Raw mean separation; sensitivity divides this distance by common sigma
  \draw[pd hairline,{Stealth[length=1.4mm]}-{Stealth[length=1.4mm]}] (axis cs:3,7.4) -- (axis cs:6,7.4);
  \node[pd label,anchor=south,align=center,text width=3.0cm] at (axis cs:4.5,7.5) {mean separation};
\end{axis}
\end{tikzpicture}
\caption{Lowering the threshold trades missed dangers for false alarms
without changing sensitivity. These Gaussian shapes use an arbitrary vertical
scale; $d'$ is their mean separation divided by their common standard deviation.
The drawn threshold is illustrative, not a fitted optimum.
[internal, Definition~\ref{def:sdt}]}
\label{fig:sdt}
\end{figure}
